\documentclass[12pt]{article}
\usepackage{amsmath, amssymb}
\usepackage[margin=1in]{geometry}
\setlength{\parskip}{6pt}
\title{QUBO Formulation: Minimum Multicut on a Tree}
\date{}
\begin{document}
\maketitle

\subsection*{Minimum Multicut on a Tree}

\paragraph{Problem Statement.}
Given a weighted tree $T = (V, E)$ with non-negative edge weights $w_e$ for each $e \in E$, and a set of terminal pairs $K = \{(s_k, t_k)\}_{k \in K}$, the Minimum Multicut on a Tree problem seeks a minimum-cost subset of edges $C \subseteq E$ such that for every $k \in K$, the terminals $s_k$ and $t_k$ reside in different connected components of the forest $T \setminus C$.

\paragraph{Decision Variables.}
For each edge $e \in E$, introduce a binary decision variable $x_e \in \{0,1\}$, where $x_e = 1$ indicates that edge $e$ is selected for removal (cut), and $x_e = 0$ indicates that edge $e$ is retained.

\paragraph{QUBO Derivation.}
\textbf{Step 1 --- Original Objective.} The goal is to minimize the total weight of the cut edges:
\begin{align}
    \min \sum_{e \in E} w_e x_e.
\end{align}

\textbf{Step 2 --- Original Constraints.} For each terminal pair $k \in K$, let $P_k \subseteq E$ denote the unique simple path between $s_k$ and $t_k$. The separation requirement is enforced by the linear inequalities:
\begin{align}
    \sum_{e \in P_k} x_e \ge 1, \quad \forall k \in K.
\end{align}

\textbf{Step 3 --- Penalty Term Conversion.} Since $x_e \in \{0,1\}$, the sum $\sum_{e \in P_k} x_e$ is integer-valued. The constraint is violated if and only if the sum equals $0$. We enforce feasibility via a quadratic penalty with weight $\lambda > 0$:
\begin{align}
    \lambda \left( \sum_{e \in P_k} x_e - 1 \right)^2.
\end{align}
Expanding the square yields:
\begin{align}
    \lambda \left( \sum_{e \in P_k} x_e^2 + \sum_{\substack{e,f \in P_k \\ e \neq f}} x_e x_f - 2 \sum_{e \in P_k} x_e + 1 \right).
\end{align}
Applying the idempotent property $x_e^2 = x_e$ for binary variables simplifies the expression to:
\begin{align}
    \lambda \left( 1 - \sum_{e \in P_k} x_e + \sum_{\substack{e,f \in P_k \\ e \neq f}} x_e x_f \right).
\end{align}
This is the general penalty contribution for constraint $k$.

\textbf{Step 4 --- Combined QUBO Hamiltonian.} Summing the objective and all penalty terms gives the total energy function $H(\mathbf{x})$:
\begin{align}
    H(\mathbf{x}) = \sum_{e \in E} w_e x_e + \sum_{k \in K} \lambda \left( 1 - \sum_{e \in P_k} x_e + \sum_{\substack{e,f \in P_k \\ e \neq f}} x_e x_f \right).
\end{align}

\textbf{Step 5 --- Q Matrix Entries and Offset.} Collecting coefficients for linear and quadratic terms over all $e, f \in E$ yields the closed-form QUBO parameters:
\begin{align}
    Q_{e,e} &= w_e - \lambda \cdot |\{k \in K : e \in P_k\}|, \quad \forall e \in E, \\
    Q_{e,f} &= \lambda \cdot |\{k \in K : \{e,f\} \subseteq P_k\}|, \quad \forall e,f \in E, e \neq f, \\
    c &= \lambda |K|.
\end{align}
The QUBO is then expressed as $H(\mathbf{x}) = \mathbf{x}^\top Q \mathbf{x} + c$.

\paragraph{Penalty Weight Justification.}
Let $W_{\max} = \max_{e \in E} w_e$. Choosing $\lambda > W_{\max}$ ensures that the penalty cost incurred by violating any single constraint strictly exceeds the maximum possible reduction in the objective function achievable by setting any subset of variables to one. More conservatively, $\lambda > \max_{k \in K} \sum_{e \in P_k} w_e$ guarantees that the penalty for infeasibility dominates the total objective value, ensuring that no infeasible configuration can achieve a lower energy than the optimal feasible one.

\paragraph{Correctness Argument.}
Minimizing $H(\mathbf{x})$ over $\{0,1\}^{|E|}$ recovers the optimal multicut solution. For any feasible assignment, all constraints are satisfied, causing every penalty term to evaluate to zero; thus $H(\mathbf{x})$ reduces exactly to the original objective function. For any infeasible assignment, at least one constraint is violated, incurring a penalty of at least $\lambda$. By the choice of $\lambda$, this penalty strictly outweighs the maximum possible decrease in the objective function that could result from flipping variables to satisfy the constraint. Consequently, the global minimum of $H(\mathbf{x})$ must occur at a feasible point, and among all feasible points, it minimizes the original objective.

\paragraph{Illustrative Example.}
Consider a tree with three edges $E = \{e_0, e_1, e_2\}$ and weights $w_0, w_1, w_2$. Suppose there is a single terminal pair ($|K|=1$) whose unique path traverses all three edges, i.e., $P_1 = \{e_0, e_1, e_2\}$. The general formulas instantiate as follows. The objective is $w_0 x_0 + w_1 x_1 + w_2 x_2$. The constraint is $x_0 + x_1 + x_2 \ge 1$. The penalty term expands to $\lambda(1 - x_0 - x_1 - x_2 + 2x_0x_1 + 2x_0x_2 + 2x_1x_2)$. The resulting QUBO Hamiltonian is:
\begin{align}
    H(\mathbf{x}) = w_0 x_0 + w_1 x_1 + w_2 x_2 + \lambda(1 - x_0 - x_1 - x_2 + 2x_0x_1 + 2x_0x_2 + 2x_1x_2).
\end{align}
The corresponding Q matrix entries and offset are:
\begin{align}
    Q_{0,0} &= w_0 - \lambda, \quad Q_{1,1} = w_1 - \lambda, \quad Q_{2,2} = w_2 - \lambda, \\
    Q_{0,1} = Q_{1,0} &= 2\lambda, \quad Q_{0,2} = Q_{2,0} = 2\lambda, \quad Q_{1,2} = Q_{2,1} = 2\lambda, \\
    c &= \lambda.
\end{align}
Assuming $w_0 \le w_1 \le w_2$, the optimal solution is $x_0 = 1, x_1 = 0, x_2 = 0$, yielding a minimum energy of $H(\mathbf{x}^*) = w_0$. This corresponds to cutting the minimum-weight edge on the path, which correctly separates the terminal pair at minimal cost.

\end{document}